% Conclusion and Future Work — ReCAHS paper (07)

\section{Conclusion and Future Work}
\label{sec:conclusion}

We set out to test a natural hypothesis --- that attention-head
importance in time-series Transformers depends on the temporal regime of
the input, and that regime-conditional head selection should therefore
beat a single static mask --- and subjected it to what we believe is the
first criterion-controlled evaluation of dynamic pruning, over
$16$ dataset$\times$horizon cells, $5$ seeds and $6$ sparsity levels.

The hypothesis survives in a sharper, redistributed form. Regime-dependent
head structure is real within every trained model, and a per-window
oracle over regime-specialist masks outperforms even the unpruned
network almost everywhere. But none of that value is accessible through
selection at inference time: once the selection criterion is held fixed,
regime routing adds nothing on temporal-patch attention, and a hierarchy
of twelve routers together with capacity, data-scaling and
distribution-shift diagnostics shows the per-window optimal mask to be
aleatoric. What routing cannot capture, combination partially can:
averaging the specialists' predictions beats every router and the best
static mask, and surpasses the unpruned model on two of four datasets ---
a diversity effect for which regime partitioning is the strongest, though
not a necessary, source.

Alongside this analysis the study yields directly usable results: joint
greedy selection replaces the fragile global-importance ranking at no
extra conceptual cost; $25\%$ of heads can be removed for free when the
mask is imposed during training; a lightweight learned detector removes
the $45\times$ STL bottleneck from any regime-aware deployment; and the
conclusions transfer to iTransformer with a precise, mechanistically
interpretable boundary --- cross-variate attention retains a small
genuine routing benefit that temporal attention does not.

More broadly, our results argue that claims about input-conditional
sparsity should be required to pass two tests they rarely face: a
selection-criterion-matched static control, and an oracle-normalized
account of how much adaptive headroom exists and how much is actually
recovered. We release per-window losses for every configuration to make
such tests cheap for future work.

Our findings reframe rather than close the adaptive-sparsity
question, and point to four directions for future work.

\paragraph{Train the specialists to be routable.}
The masks we route between are found \emph{post hoc} on a model never
trained to be routable, so their per-window advantages are accidental ---
which is exactly what the aleatoric result reflects. Learning head gates
end-to-end (a mixture-of-head-experts with a load-balancing objective, in
the spirit of conditional-computation architectures
\citep{han2021dynamic,rao2021dynamicvit}) would make routing signal
learnable by construction, and our oracle-normalized evaluation provides
the yardstick for whether trained routability closes the gap that
post-hoc routing cannot.

\paragraph{Characterize the aleatoric residual.}
The oracle gap behaves like noise to every predictor we tried, but a
formal account is missing: under what conditions on the mask-perturbation
covariance and the target noise is per-window mask identity
information-theoretically unpredictable, and how does the recoverable
fraction scale with mask diversity? The released per-window losses make
this a data-analysis problem rather than a training problem.

\paragraph{Exploit the boundary.}
Cross-variate attention retained a genuine routing benefit where temporal
attention did not. Testing whether this persists on strongly multivariate
benchmarks (ECL, Traffic, larger Weather horizons), across capacity
budgets, and for other inverted or hybrid architectures would turn a
single boundary observation into a design principle for where adaptive
head sparsity is worth building.

\paragraph{Cheaper and richer conditioning signals.}
STL regimes are one coarse conditioning variable. Change-point detectors,
learned soft regime embeddings, or exogenous covariates may induce
specialist pools with deeper \emph{and more exploitable} structure; the
diversity ablation suggests evaluating any such proposal on two axes ---
realized ensemble gain and oracle depth --- rather than on routing
accuracy alone. On the systems side, the top-$2$ ensemble's $2\times$
forward cost can likely be reduced by sharing computation between masks
that differ in only a few heads.
